\documentclass{article}
\usepackage{tikz}
\usepackage{array} % Para alinear en fila

\begin{document}
	
	\begin{center}
		\begin{tabular}{c c c}
			% Primera figura: Vectores con ángulo 0
			\begin{tikzpicture}[scale=1]
				% Dibujar ejes
				\draw[->] (-1,0) -- (3,0) node[right] {$x$};
				\draw[->] (0,-1) -- (0,3) node[above] {$y$};
				
				% Dibujar vectores
				\draw[->,thick,blue] (0,0) -- (2,1) node[above] {$\vec{a}$};
				\draw[->,thick,red] (0,0) -- (3,1.5) node[above] {$\vec{b}$};
			\end{tikzpicture} 
			&
			% Segunda figura: Vectores con ángulo 15 grados
			\begin{tikzpicture}[scale=1]
				% Dibujar ejes
				\draw[->] (-1,0) -- (3,0) node[right] {$x$};
				\draw[->] (0,-1) -- (0,3) node[above] {$y$};
				
				% Dibujar vectores con 15 grados de diferencia
				\draw[->,thick,blue] (0,0) -- (2,1) node[above] {$\vec{a}$};
				\draw[->,thick,red] (0,0) -- (2.1,1.6) node[above] {$\vec{b}$};
				
				% Dibujar semicírculo para ángulo de 15°
				\draw[thick] (0.8,0.4) arc (27:37:1);
				\node at (1,0.2) {\small $15^\circ$};
			\end{tikzpicture} 
			&
			% Tercera figura: Vectores perpendiculares (CORREGIDA)
			\begin{tikzpicture}[scale=1]
				% Dibujar ejes
				\draw[->] (-1,0) -- (3,0) node[right] {$x$};
				\draw[->] (0,-1) -- (0,3) node[above] {$y$};
				
				% Dibujar vectores perpendiculares
				\draw[->,thick,blue] (0,0) -- (2,1) node[above] {$\vec{a}$};
				\draw[->,thick,red] (0,0) -- (-1,2) node[left] {$\vec{b}$};
				
				% Dibujar ángulo de 90° entre a y b (CORRECTO)
				\draw[thick] (0.6,0.3) arc (26.57:116.57:0.7); % Arco ajustado
				\node at (0.3,1) {\small $90^\circ$}; % Posición corregida
			\end{tikzpicture} 
			\\
			{\small $\vec{b}$ un número por $\vec{a}$} & 
			{\small $\vec{b}$ casi un número por $\vec{a}$} & 
			{\small $\vec{b}$ muy distinto a $\vec{a}$ por un número}
		\end{tabular}
	\end{center}
	
\end{document}